\chapter{系统测试}

\section{测试计划}

测试工作主要验证网络连通性、区域隔离效果、访问控制策略和冗余链路可用性。测试前应准备网络拓扑图、地址规划表、设备配置文件和测试用例表。

\subsection{单元测试}

单元测试用于检查单个网络区域或单台设备的配置是否符合设计要求。测试项目包括终端地址获取、网关连通性、同区域访问和基础带宽表现。

\subsection{结合测试}

结合测试用于验证多个网络区域之间的协同运行效果。测试项目包括跨区域访问、服务器区服务访问、非法访问拦截和冗余链路切换。

\begin{table}[htbp]
  \centering
  \caption{测试用例示例}
  \label{tab:test-case}
  \begin{tabularx}{0.88\textwidth}{lllX}
    \toprule
    编号 & 测试内容 & 预期结果 & 说明 \\
    \midrule
    T01 & 教学区访问网关 & 连通 & 验证基础接入配置 \\
    T02 & 宿舍区访问办公区 & 阻断 & 验证区域隔离策略 \\
    T03 & 办公区访问服务器 & 连通 & 验证授权业务访问 \\
    T04 & 冗余链路切换 & 业务不中断或短暂恢复 & 验证核心链路可靠性 \\
    \bottomrule
  \end{tabularx}
\end{table}

\section{测试结论}

通过连通性测试、安全策略测试和冗余测试，可以确认方案是否满足校园网建设需求。若测试中发现访问控制过宽、网段划分不清或路由路径异常，应回到设计阶段进行调整。
